%% Example Resume - [YOUR_NAME]
%% A ONE-PAGE, role-targeted document. Use it when the posting wants a resume,
%% when you are early-career, or when you want a tighter read than the CV.
%% For the comprehensive document (publications, honors, full education
%% history) use cv/main_example.tex instead - that one is the 2-page CV.
%%
%% What actually differs from the CV, and why:
%%   - Hard 1-page budget (the CV is 2). Achievements get cut, not type size.
%%   - A smaller name block, to buy vertical room.
%%   - No Publications / Honors / References sections. Those are CV-only; on a
%%     resume they cost a third of the page and read as padding.
%%   - Skills as one dense line instead of a bulleted competencies list.
%%   - Projects carries its own section (the score-carrying section for
%%     students and freshers) - see 05-cv-templates.md's stage ordering.
%%
%% The slots below are sized so a FILLED resume lands at roughly 80% of the
%% page, leaving headroom for real achievement text to run 1-2 lines each
%% without breaking the 1-page budget. Delete a slot you do not need; do not
%% shrink the type or the geometry scale to make something fit.
%%
%% SETUP: Run /setup to populate this with your actual information.
%% Compile with: lualatex resume_example.tex
%% (pdflatex often fails on modern MiKTeX with fontawesome5 font-expansion errors.)

\documentclass[10pt,a4paper,sans]{moderncv}
\moderncvstyle{banking}
\moderncvcolor{blue}

% Same name/section colouring fix as the CV (moderncvstylebanking.sty's
% \colorlet copies, not aliases, freeze the name colours before
% \moderncvcolor runs). Smaller than the CV's 34pt to save vertical space.
\renewcommand*{\namefont}{\fontsize{26}{28}\bfseries\upshape}
\colorlet{firstnamecolor}{color1}
\colorlet{lastnamecolor}{color1}
\colorlet{namecolor}{color1}
\renewcommand*{\sectionstyle}[1]{{\sectionfont\color{color1}#1}}

\usepackage[utf8]{inputenc}
% moderncv loads hyperref itself in an \AtEndPreamble hook, so \hypersetup
% must go in an \AtEndPreamble of our own (see the note in 05-cv-templates.md).
\AtEndPreamble{\hypersetup{
    colorlinks=true,
    linkcolor=blue,
    filecolor=magenta,
    urlcolor=blue,
    pdftitle={[YOUR_NAME] - Resume},
    pdfpagemode=UseNone,
}}
% Same scale as the CV, and deliberately so: geometry's `scale` sets the TEXT
% BLOCK WIDTH, not the type size (the size is the documentclass font above,
% 10pt here vs the CV's 11pt). A lower value narrows the column, which is
% strictly bad for a one-pager: the same content needs MORE lines, and the
% measure gets cramped and hard to scan. It saves nothing. This document
% reaches 1 page on content alone, at full width.
\usepackage[scale=0.80]{geometry}

% personal data
\name{[First]}{[Last]}
% India convention: city + state only, no street address or PIN.
\address{[City, State]}{}{}
% The placeholder must stay underscore-free: LaTeX reads a bare `_` as a
% subscript and dies with "Missing $ inserted" (F31 in tests/test_latex_guidance.py).
\phone[mobile]{+91 XXXXX XXXXX}
\email{[your.email@example.com]}
\extrainfo{\href{[https://linkedin.com/in/your-profile]}{LinkedIn}, \href{[https://github.com/your-username]}{GitHub}}

\begin{document}

\makecvtitle

% ============================================================
%     PROFILE STATEMENT
% ============================================================
% Two lines, not 3-5: on a one-pager this is a claim, not a biography.
% Lead with the single strongest match to THIS posting.

\vspace{4pt}
\small{[One or two lines naming the role you are targeting and the one thing that makes you a fit for it. Example: "Machine-learning engineer with X years building production NLP systems, from retrieval to evaluation, in Python and PyTorch."]}

% ============================================================
%     SKILLS
% ============================================================
% One dense line beats a bulleted list here: it costs a third of the space
% and the ATS keyword scan reads it the same way.

\section{Skills}
\vspace{1pt}
\small{\textbf{Languages:} [Language 1, Language 2]. \textbf{Frameworks:} [Framework 1, Framework 2]. \textbf{Tools:} [Tool 1, Tool 2]. \textbf{Domain:} [Domain or vertical].}

% ============================================================
%     EXPERIENCE
% ============================================================
% Three slots: two roles plus one internship. For students and freshers the
% internship is the strongest single signal on the page - if you have done
% one, that is the slot to keep and a part-time role to drop. If you have
% more than three, the CV (cv/main_example.tex) is the document that holds
% them.

\section{Experience}
\vspace{2pt}
\begin{itemize}

% --- Most Recent Role ---
\item{\cventry{[YYYY-Present]}{[Job Title]}{[Company]}{[City, Country]}{}{\vspace{1pt}
\begin{itemize}
    \item {[Achievement with a number in it - strongest first]}
    \item {[Second achievement]}
\end{itemize}}}

\vspace{2pt}

% --- Previous Role ---
\item{\cventry{[YYYY-YYYY]}{[Job Title]}{[Company]}{[City, Country]}{}{\vspace{1pt}
\begin{itemize}
    \item {[Achievement]}
    \item {[Achievement]}
\end{itemize}}}

\vspace{2pt}

% --- Internship (students / freshers) ---
% For an experienced candidate, replace this slot with a third role or
% delete it - do not leave it as a placeholder.
\item{\cventry{[YYYY]}{[Internship Role]}{[Company]}{[City, Country]}{}{\vspace{1pt}
[What you built or did, and the outcome. One or two lines.]}}

\end{itemize}

% ============================================================
%     PROJECTS
% ============================================================
% The score-carrying section for students and freshers - keep it above
% Education. For an experienced candidate, cut it if the page is tight and
% the role does not ask for it.

\section{Projects}
\vspace{1pt}
\begin{itemize}

\item{\cventry{[YYYY]}{[Project Name]}{[Stack or domain]}{}{}{\vspace{1pt}
[What it does, what you built, and the outcome. Link it if the template
allows.]}}

\vspace{2pt}

\item{\cventry{[YYYY]}{[Project Name]}{[Stack or domain]}{}{}{\vspace{1pt}
[One line - outcome first.]}}

\vspace{2pt}

\item{\cventry{[YYYY]}{[Project Name]}{[Stack or domain]}{}{}{\vspace{1pt}
[One line - outcome first.]}}

\end{itemize}

% ============================================================
%     CERTIFICATIONS
% ============================================================
% Online course certificates, competitive-programming ranks, hackathon
% wins. This is NOT the CV's Honors section: one line per certificate, no
% dates, no issuer boilerplate. Delete the whole section if you have none -
% an empty "Certifications" heading reads worse than its absence.

\section{Certifications}
\vspace{1pt}
\begin{itemize}

\item{[Certification or rank] - [Issuer], [YYYY]}
\item{[Certification or rank] - [Issuer], [YYYY]}

\end{itemize}

% ============================================================
%     EDUCATION
% ============================================================
% One or two entries. Drop thesis detail: it belongs in the CV.

\section{Education}
\vspace{1pt}
\begin{itemize}

\item{\cventry{[YYYY-YYYY]}{[Degree] in [Field]}{[Institution]}{[City, Country]}{}{\vspace{1pt}
[CGPA if strong. One line of relevant coursework or focus, no more.]}}

\end{itemize}

\end{document}
